\documentclass[11pt]{article}
\usepackage[margin=1in]{geometry}
\usepackage{amsmath,hyperref}
\title{Guarding a Role Transition While Its Earlier Work Is Still Owed}
\date{}
\begin{document}
\maketitle

\paragraph{Question and scope.}
The admission criterion in \href{https://github.com/eternal-november/eos-research-programmes/pull/353}{\#353} installs the new gate holder in both histories and rejects an optional request in one of them. This companion makes the disputed work mandatory and asks whether the proposed promotion itself can occur. In the modeled cut it can occur one tick earlier when an old closure uses an ordinary inspector, even though the proposed new holder, credentials, health report and requested promotion are the same.

\paragraph{Roles, permissions and live work.}
Civil maintenance employees $A,B,C,D$ each supply at most one named worker-time slot per tick. A continuously needed relay consumes one eligible slot per tick. Gate control is exclusive: a resource-side interlock recognizes one holder and epoch, and only that pair can admit new work or serve a holder-bound ticket. A tote closure has its own named permit, whose authorized executor is checked independently of gate control. $A$ can close an intact pocket; $D$ can close either pocket. $C,D$ can relay in ticks 2 and 3; $B$ can relay at tick 2 only if the missing health report is delayed rather than a real radio failure. $B$'s gate actuator is separately certified operable through tick 2 in both completions. The decision uses the certified pocket state and this actuator certificate, not knowledge of which radio world occurred. $C$'s appointment, authority credentials and capacity are the same in both pocket states.

At tick 0 holder $A$ admits the first tote, due for closure at 2, while $D$ relays. At tick 1 a planned transfer installs $B$ before $B$ admits the second tote, due for closure by $D$ at 3, and an already owed gate ticket due at 2. The ticket's service is not another optional admission: it consumes the gate holder's worker slot and cannot be dropped. $C$ and $D$ overlap relay coverage during the planned handoff without sharing gate permission. Both totes and the gate ticket are live at the tick-2 proposed $B\to C$ transition. The modeled schedule must preserve the four relay ticks and all three due tasks.

\paragraph{Interlock rule and exact criterion.}
The modeled authority service is the only writer and its epoch update and ticket assignment take effect atomically. A proposed successor must present the authorized acceptance record for the current epoch and a complete, current named ledger of live duties. The check takes place before allocating any work at the cut. A closure with its own permit may stay with an eligible operations employee; it does not become authorization to use the gate. A still-pending gate ticket must move, with its job identity, deadline, successor acceptance and authority, in the same guarded commit as the epoch change. If that operation cannot be certified, either the incumbent must perform the ticket before the promotion or the proposed promotion cannot occur within the stated service laws. A rewrite of the role label cannot cancel a promised task.

More explicitly, for a proposed holder $H$, retain the independently permitted closure edges, put the gate ticket on $H_2$ while still owed, and retain only relay edges legal in both health completions. For the stated finite, fixed-schedule envelope there is a legal immediate promotion if and only if authorization and ledger conditions hold and those mandatory jobs have a disjoint slot assignment. For necessity, every successful plan must spend its gate ticket under the current holder and use only common legal edges: no later observation arrives before the relevant work. Distinct duties cannot spend the same slot. The authorization and ledger conditions are legal preconditions, not inferences from capability. For sufficiency, choose the disjoint assignment before the cut, atomically install the accepted successor and ticket, and run the installed guarded operations. Every caller presents its pair; no stale pair may spend the new gate epoch. Separate closure permits are preserved and their execution uses the selected eligible employees. A physical finish settles a duty, not just the receipt of an installation notice. No maximum detection-delay result or implementation of a replicated interlock is implied.

With an intact pocket and correct acceptance evidence the early promotion has the schedule $A_2$ (first closure), $C_2$ (gate ticket), $D_2$ (relay), $D_3$ (second closure), $C_3$ (relay). With a damaged pocket, those first three tasks have only $C_2,D_2$ if $C$ takes control now. The Hall deficit is one, and the failure persists even with both authorization signatures. If the incumbent instead retains the gate for tick 2, $B_2$ services the ticket, $D_2$ closes the damaged pocket and $C_2$ relays. At tick 3 the gate can pass to $C$ once that binding ticket has actually settled, with the second tote still owed on a separate permit before it closes later that tick. Invalid evidence defers rather than inventing a valid transfer. An older queued $A$ request is rejected in both scenes, and the same old $B$ pair is rejected after either effective cut. The interlock prevents double admissions but does not by itself establish feasibility of the promotion.

\paragraph{Comparison and replay.}
\href{https://robotics.usc.edu/publications/347/}{Gerkey and Matari\'c} and \href{https://arxiv.org/abs/2105.05586}{Notomista et al.} respectively provide assignment and resilient capability baselines, not this model's named legacy permission or effective-cut rules. Given those laws, factored direct slot assignment and the \#353 exact spare-profile admission calculation agree; neither creates authority. The extra guarantee is that the same atomic epoch operation preserves previously owed service, or an exact obstruction to installing it immediately. The replay independently checks feasibility by slot permutations, traces 16 reached cases across pocket, actual hidden health and evidence, and checks continuous relay, physical completion and stale-message rejection. All results concern the stated warehouse simulation.
\end{document}
